\section{Political Economy Discussion: Competing Paradigms, Empirical Causality, and Peripheral Solvency}
\label{sec:discussion_conclusion}

\subsection{Competing Paradigms in Economic Thought: Evaluating Research Programmes on the Unidad Popular}
\label{sec:competing_paradigms}

\paragraph{Rival Worldviews on the Unidad Popular.}
In the history of economic thought, theoretical advance rarely proceeds through neutral empirical accumulation. As \citet{Roncaglia2006} observes, analytical progress unfolds through an enduring contest among rival paradigms rooted in conflicting conceptual frameworks and social visions. For more than half a century, mainstream economics has interpreted the breakdown of Chile's Unidad Popular (UP, 1970--1973) through the lens of ``macroeconomic populism'' \citep{DornbuschEdwards1990,Edwards2024}. In this narrative, unforced fiscal expansion and reckless money printing triggered an inevitable cycle of hyperinflation and supply collapse. Yet this orthodox consensus analyzes peripheral capitalism as a closed monetary system, abstracting from international currency hierarchies, terms-of-trade collapses, and external financial constraints. Evaluated through \citeauthor{Lakatos1978}'s (\citeyear{Lakatos1978}) methodology of scientific research programmes, the orthodox interpretation exhibits the classic symptoms of theoretical stagnation, while dependency theory as a research programme \citep{Kvangraven2021} provides a progressive framework capable of explaining the empirical causal record.

\paragraph{The Monetarist Paradigm and Its Explanatory Anomalies.}
A research programme degenerates when its theoretical core produces no novel predictions and survives primarily by introducing ad hoc adjustments to accommodate anomalies \citep{Lakatos1978}. The orthodox monetarist programme rests on three foundational axioms: the exogenous determination of the money supply, the direct quantity-theoretic pass-through from money to prices, and long-run monetary neutrality under full-capacity market clearing. Confronted with the historical record of 1971---when base money expanded by over 100\% while inflation slowed from 35\% to 22\% and manufacturing output expanded by 12.1\%---orthodoxy relies on narrative rationalizations. It dismisses the cancellation of \$294~million in foreign commercial credit lines, the U.S. Eximbank embargo, and falling copper prices as secondary disturbances \citep{Edwards2024}. In addition, its reliance on market-clearing assumptions leads it to diagnose all retail shortages as ``repressed inflation'' caused by price controls, ignoring organized employer distribution lockouts and the fact that manufacturing capacity utilization reached 98.4\% alongside falling net fixed investment ($\Delta K = -2.3\%$). Over fifty years, this paradigm has produced few novel insights, repeating a monocausal narrative that overlooks the timing, direction, and non-linearities of peripheral monetary instability.

\paragraph{Dependency Theory as a Progressive Research Programme.}
In contrast, treating dependency theory as a research programme \citep{Prebisch1950,Marini1973,DosSantos1970,Kvangraven2021} generates excess empirical content. Rather than modeling peripheral economies as scaled-down versions of advanced capitalist nations, dependency theory anchors macroeconomic outcomes in structural asymmetries: external financial subordination, specialized export structures, and an unequal international currency hierarchy \citep{Vasudevan2009,Basu2022}. This chapter advances this tradition by formalizing central bank balance-sheet solvency as an endogenous state variable ($S_t \equiv F_t / H_t$). This analytical framework anticipates empirical patterns that orthodoxy cannot explain: the state-dependent breakdown of monetary transmission, the bidirectional predictive interaction with dominant reverse accommodation from price shocks to base money growth, and the loss of the external solvency buffer once foreign exchange reserves are depleted.

\subsection{Empirical Causality and Dynamic Transmission: Precedence and Non-Linear Multipliers}
\label{sec:causality_transmission}

\paragraph{Directional Precedence and the Rejection of Money Exogeneity.}
The econometric findings in Section~\ref{sec:historical_empirical_results} operate as an empirical test between these competing frameworks. Standard monetarist accounts treat central bank emission as an unconstrained, discretionary policy impulse that initiates inflation. The Granger causality tests within stationary vector autoregressions in Section~\ref{sec:granger_tournament} evaluate this assumption. Temporal precedence exhibits bidirectional predictive interaction with pronounced reverse accommodation running from consumer price inflation to base money growth ($\pi \to g_H$), with $F$-statistics four times larger than forward money growth and persisting across all tested lag horizons ($k=1,\dots,4$), while monthly central bank solvency ratio growth systematically predicts base money expansion ($g_{\text{SolvR\_H}} \to g_H$). In parallel, the constraint-based PCMCI causal discovery algorithm in Appendix~\ref{app:pcmci_ee1} retains directed lagged conditional links connecting gross reserve depletion to subsequent base money expansion and relative price distortions, though PCMCI does not retain a direct conditional $g^H \to \pi$ edge under its conditioning set ($q > 0.10$). While this method-specific finding highlights conditioning sensitivity across estimators, the VAR and GIRF estimations retain a statistically significant forward response alongside dominant reverse accommodation. Rather than acting primarily as an autonomous driver of inflation, the Central Bank functioned as an endogenous accommodation valve, issuing credit to sustain payrolls and prevent production halts as external constraints tightened.

\paragraph{Dynamic Multiplier Evaluation: Reverse Accommodation Dominance.}
The non-linear Threshold Vector Autoregression (TVAR) and Generalized Impulse Response Function (GIRF) estimates in Section~\ref{sec:tvar_girf_multipliers} provide quantitative evidence against the quantity-theoretic narrative. In this dynamic evaluation, reverse defensive accommodation ($\pi_m \to g^H$) is more persistent and cumulatively larger over the reported horizon than forward monetary transmission ($g^H \to \pi_m$). As shown in Figure~\ref{fig:tvar_girf_tournament}, Table~\ref{tab:app_d01_girf_differences}, and Table~\ref{tab:app_d02_girf_metrics}, a one-standard-deviation base money shock produces only a modest price response: domestic inflation rises by 0.97 percentage points at month 1, loses statistical significance after 9 months, and yields a cumulative 12-month multiplier of 3.55 percentage points. In contrast, a one-standard-deviation price shock triggers an immediate, sustained monetary injection: Central Bank base money expands by 1.72 percentage points at month 1 and remains statistically significant for 10 consecutive months ($h=1$ to $10$). This generates a cumulative 12-month base money emission of 6.54 percentage points. The point-wise multiplier difference test ($\Delta(h)$) confirms that reverse accommodation significantly exceeds forward transmission across 9 consecutive months ($h=2$ to $10$). When enterprise costs escalated, the monetary authority was compelled to accommodate credit demand.

\paragraph{State-Dependent Transmission: External Buffers and Balance-Sheet Insolvency.}
The TVAR estimates demonstrate that monetary transmission is state-dependent on central bank reserve solvency. In the non-adverse state, international reserves provide an effective external buffer: world price shocks ($g^{gold} \to \pi_m$) exert zero statistically significant effect on domestic consumer prices over a 12-month horizon (Panel~A of Figure~\ref{fig:tvar_girf_gold_shield}). The central bank absorbs import cost pressures by drawing on foreign exchange liquidity. Once reserve depletion breaches the estimated structural threshold in the flow growth gap ($\Delta s_{t-1} \le -4.615\%$, entering the adverse solvency-growth state), this external buffer dissolves. In the adverse state (reflecting conditions of central-bank insolvency), the identical world price shock unleashes an unbroken 10-month inflationary wave ($h=3$ to $12$), peaking at 0.98 percentage points at month 3 and generating 6.54 percentage points of cumulative inflation. The endogenous regime-switching probabilities in Table~\ref{tab:app_d03_regime_switching} indicate that this adverse state exhibits strong persistence: export price windfalls produce negligible shifts in the simulated probability of exiting the adverse state ($\Delta P \approx 0$). Under these adverse conditions, unbuffered international price shocks generate persistent domestic inflation, inducing defensive monetization to sustain production and employment.

\subsection{The International Currency Hierarchy, Hegemonic Contingency, and Central Bank Agency}
\label{sec:currency_hierarchy_agency}

\paragraph{The Peripheral Paradox: Domestic Credit Tokens versus International Reserve Money.}
These dynamic responses reflect what I define as the \emph{Peripheral Paradox} (Section~\ref{sec:money_endogeneity_paradox}). The central bank balance sheet in dependent economies embodies two distinct monetary forms. On the liability side, domestic base money ($H_t$) consists of credit tokens issued under national sovereign authority \citep{Basu2022}. On the asset side, international reserves ($IR_t$) represent claims on world money that the peripheral state cannot emit \citep{Vasudevan2009,AlamiEtAl2022}. The peripheral state can expand domestic currency liabilities, but it cannot manufacture international reserve assets. When foreign trade credit lines are severed and export receipts are blocked, the domestic monetary circuit decouples from international settlement. The central bank faces an intractable institutional dilemma: either it withholds liquidity from enterprises facing import price explosions, triggering immediate insolvency and layoffs; or it acts as lender of last resort, expanding credit that rapidly depreciates the currency in the absence of foreign exchange backing.

\paragraph{Hegemonic Contingency, Central Bank Agency, and the Bretton Woods Environment.}
The Chilean transition unfolded against the breakdown of the post-war international monetary order. As \citet{Zoeller2019} demonstrates, President Nixon's decision to close the gold window on August 15, 1971, was not an inevitable structural transition, but a tactical contingency plan deployed by the United States to preserve domestic policy autonomy and manage balance-of-payments pressures. It is vital, however, not to conflate this global shift with the domestic crisis in Chile. The dissolution of Bretton Woods altered the international macroeconomic environment, exporting global inflation and exchange-rate instability. Yet the Central Bank of Chile (BCCh) faced a targeted, bilateral financial blockade: private commercial banks canceled \$220~million in short-term credit lines, while the U.S. Eximbank and multilateral lenders halted disbursements \citep{DeVylder1976,PetrasMorley1975}. In this context, the BCCh exercised institutional agency. Rather than passively reflecting global turbulence or imposing an orthodox deflationary adjustment on employment, the central bank chose to accommodate enterprise credit demand and finance essential food imports, consciously absorbing the external squeeze on its own balance sheet.

\paragraph{Institutional Limits of Currency Sovereignty: A Critique of Naive Chartalism.}
This historical experience offers a clear critique of naive extensions of Modern Monetary Theory (MMT) to the Global South. Neo-chartalist claims that sovereign currency issuers face no financial budget constraints reflect the institutional realities of core reserve-issuing economies \citep{AlamiEtAl2022,Vasudevan2009}. In an asymmetric currency hierarchy, formal legal authority over domestic fiat money does not grant command over the universal means of international settlement. A peripheral government cannot overcome balance-of-payments bottlenecks by issuing domestic credit. Without foreign exchange reserves to purchase non-substitutable capital equipment, spare parts, and agricultural inputs, domestic liquidity injections lead to currency collapse and imported cost-push inflation. Progressive development policy in the periphery cannot bypass the structural constraint of external solvency.

\subsection{Synthesis, Methodological Scope, and Concluding Remarks}
\label{sec:synthesis_horizons}

\paragraph{Evidentiary Limits and Econometric Caveats.}
A rigorous appraisal must acknowledge the evidentiary limits and empirical scope of macro time-series econometrics. The empirical methods deployed in this chapter---stationary vector-autoregressive Granger causality tests, PCMCI causal discovery, and non-linear threshold vector autoregressions---establish directional precedence and state-dependent conditional responses. They identify systematic empirical patterns across historical regimes, but they do not constitute micro-level counterfactual identification in the experimental sense. Macro aggregates remain subject to the limitations of unobserved common shocks and structural breaks across historical periods. These econometric findings provide evidence consistent with balance-sheet constraints and endogenous accommodation, but they should be viewed as opening, rather than settling, the empirical inquiry into peripheral monetary institutions.

\paragraph{Analytical Findings and Directions for Future Research.}
By grounding the macroeconomic history of the Unidad Popular in an open-economy, balance-sheet framework, this chapter demonstrates the progressive capacity of dependency theory as a research programme \citep{Kvangraven2021}. The analysis accounts for anomalies that disrupt orthodox narratives, replacing static monetarist assumptions with an empirical model of reserve solvency. Rather than offering definitive closures, this framework points toward several avenues for future research. One important frontier involves examining the meso-scale institutional mechanisms of credit allocation, investigating how credit flows were mediated between the central bank, commercial banks, and state enterprises during periods of balance-of-payments distress. Another promising avenue is comparative: exploring whether the solvency threshold and reserve buffering dynamics identified in Chile operate across other peripheral economies confronting external debt and currency crises.

\paragraph{Institutional Implications: Foreign-Exchange Management and Structural Constraints.}
The historical experience of the Unidad Popular therefore points beyond the problem of macroeconomic management narrowly understood. Foreign-exchange budgeting, controls over capital movements, and the regulation of external financial flows can widen the policy space available to peripheral governments and reduce their exposure to destabilizing international liquidity cycles \citep{FfrenchDavis2018}. Yet greater room for maneuver within the balance-of-payments constraint does not by itself transform the relations that reproduce that constraint. A peripheral state may command the domestic means of payment while remaining unable to produce the world money required to reproduce an import-dependent productive structure. Foreign exchange is therefore more than an object of prudent macroeconomic administration: it is a material condition of accumulation and an object of conflict over the control of export revenues, external credit, imported means of production, and the surplus through which productive capacity is reproduced. Macroprudential autonomy can mitigate the vulnerability generated by dependent accumulation, but it cannot by itself overcome it. The problem posed by peripheral transformation is consequently not only how to manage the balance-of-payments constraint, but how to alter the domestic and international relations through which that constraint is continually reproduced.
